\newcommand{\figuraPerfilLosaIsotermica}{%
\begin{figure}[H]
    \centering
    \begin{tikzpicture}[>=Latex,font=\small]
        \begin{axis}[
            width=11.8cm,
            height=7.2cm,
            xmin=-4,xmax=4,
            ymin=0,ymax=1.12,
            axis lines=middle,
            xlabel={$Z$},
            ylabel={$\rho/\rho_0$},
            xtick={-4,-3,-2,-1,0,1,2,3,4},
            ytick={0,0.2,0.4,0.6,0.8,1},
            grid=major,
            grid style={gray!18},
            title={Perfil de una losa isotérmica autogravitante},
            samples=180,
            clip=false]

            \addplot[very thick,blue!70!black,domain=-4:4]
                {1/(cosh(x)^2)};

            \draw[densely dashed,red!70!black]
                (axis cs:0,0)--(axis cs:0,1.05);
            \node[red!70!black,anchor=south east]
                at (axis cs:-0.10,1.01) {plano de simetría};

            \draw[->,ultra thick,green!45!black]
                (axis cs:-3.2,0.36)--(axis cs:-2.35,0.36)
                node[midway,above] {$\vec{g}$};
            \draw[->,ultra thick,green!45!black]
                (axis cs:3.2,0.36)--(axis cs:2.35,0.36)
                node[midway,above] {$\vec{g}$};

            \node[blue!70!black,anchor=north east]
                at (axis cs:3.75,0.91) {perfil de densidad};
        \end{axis}
    \end{tikzpicture}
    \caption{La densidad de la losa alcanza su máximo $\rho_0$ en el plano
    de simetría y decrece como $\operatorname{sech}^2 Z$. El campo
    gravitacional apunta desde ambos lados hacia el plano central.}
    \label{fig:perfil-losa-isotermica}
\end{figure}%
}